\section{Аналіз помилок застосування асиметричної криптографії}

\subsection{Ситуація 1. Неправильне використання padding у RSA}

\textbf{Умова:} RSA використовується без рекомендованої схеми padding або із застарілою схемою,
непридатною для відповідного сценарію.

\textbf{Помилка:} Використання сирого (textbook) RSA або застарілої схеми PKCS\#1 v1.5 замість сучасних схем
доповнення RSA-OAEP / RSA-PSS.

\textbf{Можливі наслідки:} Вразливість системи до атак типу Padding Oracle (атака Блейхенбахера), можливість
підробки цифрових підписів, детермінованість шифротексту та ризик розшифрування повідомлень малої
довжини через добування кореня.

\textbf{Спосіб усунення:} Обов'язкове переведення системи на стандарт RSA-OAEP для шифрування та
RSA-PSS для цифрових підписів відповідно до вимог RFC 8017~\cite{RFC8017}.

\subsection{Ситуація 2. Повторне використання nonce в ECDSA}

\textbf{Умова:} Під час створення двох різних ECDSA-підписів було повторно використано однаковий nonce.

\textbf{Помилка:} Використання однакового або передбачуваного тимчасового параметра $k$ (nonce) під час
створення кількох підписів ECDSA.

\textbf{Можливі наслідки:} Повне математичне виведення та розкриття приватного ключа підписувача на
основі двох відомих підписів і відповідних повідомлень. Повна компрометація системи.

\textbf{Спосіб усунення:} Застосування алгоритмів із детермінованою генерацією $k$ на основі хешу
повідомлення та приватного ключа (стандарт RFC 6979) або перехід на алгоритм EdDSA
(Ed25519)~\cite{RFC6979,RFC8032}.

\subsection{Ситуація 3. Передбачуваний генератор випадкових чисел}

\textbf{Умова:} Криптографічні ключі генеруються з використанням передбачуваного генератора випадкових
чисел.

\textbf{Помилка:} Застосування некриптографічних генераторів псевдовипадкових чисел (наприклад,
\texttt{random.randint()} у Python) замість криптографічно стійких CSPRNG.

\textbf{Можливі наслідки:} Значне звуження простору ключового матеріалу, можливість відтворення
послідовності генератора зловмисником та швидкий повний перебір (brute-force) створених ключів.

\textbf{Спосіб усунення:} Використання виключно криптографічно стійких генераторів випадкових чисел
(CSPRNG), підкріплених джерелами ентропії операційної системи (\texttt{os.urandom()}, \texttt{secrets},
системні виклики \texttt{getrandom()})~\cite{NIST80057}.

\subsection{Ситуація 4. Відсутність перевірки відкритого ключа}

\textbf{Умова:} Система отримує відкритий ключ від іншої сторони та використовує його без перевірки його
належності очікуваному учаснику взаємодії.

\textbf{Помилка:} Нехтування автентифікацією відкритого ключа та відсутність перевірки цифрових
сертифікатів партнерів по взаємодії.

\textbf{Можливі наслідки:} Успішна реалізація атаки <<людина посередині>> (Man-in-the-Middle).
Зловмисник підміняє відкритий ключ власним, що дозволяє йому перехоплювати, розшифровувати та
модифікувати весь трафік.

\textbf{Спосіб усунення:} Впровадження інфраструктури відкритих ключів (PKI), перевірка ланцюжків
сертифікатів X.509 або використання механізмів фіксації ключів (Key Pinning)~\cite{NIST80057}.

\subsection{Ситуація 5. Повторне використання ключової пари}

\textbf{Умова:} Одна й та сама пара ключів використовується одночасно для різних криптографічних
призначений без аналізу сумісності таких способів використання.

\textbf{Помилка:} Мультифункціональне використання однієї пари ключів (наприклад, одночасне
застосування однієї RSA-пари для шифрування даних і для накладання цифрових підписів).

\textbf{Можливі наслідки:} Перехресні атаки між різними протоколами. Зловмисник може змусити систему
підписати зашифроване повідомлення, перетворюючи сервіс підпису на оракул розшифрування.

\textbf{Спосіб усунення:} Суворе дотримання принципу розділення ключів (Key Separation) шляхом випуску
окремих пар ключів для кожної криптографічної операції та контролю прапорців призначення в сертифікатах
(\texttt{KeyUsage})~\cite{NIST80057}.

\subsection{Загальний висновок щодо помилок застосування}

Наведені ситуації демонструють, що математична стійкість криптографічного алгоритму сама по собі
негарантує безпечності інформаційної системи. Більшість критичних вразливостей виникає не через злам
математичної основи шифрів, а внаслідок помилок їх реалізації та архітектурного застосування: використання
передбачуваних джерел випадковості, некоректного вибору схем доповнення (padding), повторного
використання тимчасових параметрів (nonce), відсутності автентифікації відкритих ключів та порушення
принципу розділення ключів. Для забезпечення реальної безпеки необхідно чітко дотримуватись протоколів
та вимог міжнародних криптографічних стандартів.
